\documentclass{article}
\usepackage{amsmath,amssymb,amsthm}
\usepackage{algorithm}
\usepackage{algpseudocode}
\usepackage{fullpage}
\newtheorem{theorem}{Theorem}
\newtheorem{lemma}{Lemma}
\newtheorem{assumption}{Assumption}
\newcommand{\E}{\mathbb{E}}
\newcommand{\norm}[1]{\left\|#1\right\|}
\newcommand{\R}{\mathbb{R}}

\begin{document}

\section*{RMSProp (Running Average of Squared Gradients)}

\section*{Mathematical Formulation}
Let $f:\R^{d}\rightarrow\R$ be differentiable.  
At iteration $t\ge 1$ we have a stochastic gradient $g_{t}\in\R^{d}$ such that  
\[
g_{t}= \nabla f(w_{t};\xi_{t}),\qquad \xi_{t}\sim\mathcal{D}
\]
where $\xi_{t}$ denotes the random data sample.  

\begin{align}
v_{t} &= \beta\,v_{t-1}+(1-\beta)\,g_{t}\odot g_{t}, \qquad v_{0}=0 \label{eq:vt}\\[4pt]
w_{t+1} &= w_{t}-\frac{\alpha}{\sqrt{v_{t}+\varepsilon}}\odot g_{t}, \qquad w_{0}\in\R^{d} \label{eq:wt}
\end{align}
Here $\beta\in[0,1)$ is the decay rate, $\alpha>0$ the base learning rate, $\varepsilon>0$ a small constant to avoid division by zero, and $\odot$ denotes element‑wise product.  
The denominator $\sqrt{v_{t}+\varepsilon}$ is also taken element‑wise.

\section*{Pseudocode}
\begin{algorithm}
\caption{RMSProp}
\begin{algorithmic}[1]
\State \textbf{Input:} $w_{0}\in\R^{d}$, $\alpha>0$, $\beta\in[0,1)$, $\varepsilon>0$, $T$.
\State $v_{0}\gets 0\in\R^{d}$
\For{$t=0$ \textbf{to} $T-1$}
    \State Sample $\xi_{t}\sim\mathcal{D}$ \Comment{draw a minibatch}
    \State $g_{t}\gets \nabla f(w_{t};\xi_{t})$
    \State $v_{t+1}\gets \beta v_{t}+(1-\beta) g_{t}\odot g_{t}$ \Comment{Eq.~\eqref{eq:vt}}
    \State $w_{t+1}\gets w_{t}-\dfrac{\alpha}{\sqrt{v_{t+1}+\varepsilon}}\odot g_{t}$ \Comment{Eq.~\eqref{eq:wt}}
\EndFor
\State \textbf{Return} $w_{T}$
\end{algorithmic}
\end{algorithm}

\section*{Dry Run Example}
Consider the scalar function $f(w)=(w-3)^{2}$, thus $\nabla f(w)=2(w-3)$.  
Set $w_{0}=0$, $\alpha=0.1$, $\beta=0.9$, $\varepsilon=10^{-8}$.  
We perform three iterations.

\begin{center}
\begin{tabular}{c|c|c|c|c}
$t$ & $w_{t}$ & $g_{t}=2(w_{t}-3)$ & $v_{t}$ & $w_{t+1}=w_{t}-\dfrac{\alpha}{\sqrt{v_{t}+\varepsilon}}g_{t}$\\\hline
0 & $0$ & $-6$ & $v_{1}=0.1\cdot 36=3.6$ & $w_{1}=0-\dfrac{0.1}{\sqrt{3.6}}(-6)=0+0.1\cdot\frac{6}{1.897}=0.316$\\[4pt]
1 & $0.316$ & $-5.368$ & $v_{2}=0.9\cdot3.6+0.1\cdot(5.368)^{2}=3.24+0.1\cdot28.82=6.122$ &
$w_{2}=0.316-\dfrac{0.1}{\sqrt{6.122}}(-5.368)=0.316+0.1\cdot\frac{5.368}{2.474}=0.736$\\[4pt]
2 & $0.736$ & $-4.528$ & $v_{3}=0.9\cdot6.122+0.1\cdot(4.528)^{2}=5.5098+0.1\cdot20.50=7.5598$ &
$w_{3}=0.736-\dfrac{0.1}{\sqrt{7.5598}}(-4.528)=0.736+0.1\cdot\frac{4.528}{2.749}=0.901$\\
\end{tabular}
\end{center}
The objective values are $f(w_{0})=9$, $f(w_{1})\approx7.78$, $f(w_{2})\approx5.92$, $f(w_{3})\approx4.90$, showing monotone descent.

\newpage
\section*{Assumptions}
\begin{assumption}[Smoothness]\label{ass:smooth}
$f$ is $L$‑smooth: for all $x,y\in\R^{d}$,
\[
f(y)\le f(x)+\langle\nabla f(x),y-x\rangle+\frac{L}{2}\norm{y-x}^{2}.
\]
\end{assumption}
\begin{assumption}[Unbiased Gradient with Bounded Variance]\label{ass:unbiased}
For every $t$, $\E[g_{t}\mid w_{t}]=\nabla f(w_{t})$ and 
\[
\E\big[\norm{g_{t}-\nabla f(w_{t})}^{2}\mid w_{t}\big]\le\sigma^{2}.
\]
\end{assumption}
\begin{assumption}[Bounded Second Moment]\label{ass:bound}
There exists $G>0$ such that $\E[\norm{g_{t}}^{2}]\le G^{2}$ for all $t$.
\end{assumption}
\begin{assumption}[Hyper‑parameter Range]\label{ass:hyper}
$\beta\in[0,1)$, $\alpha>0$, $\varepsilon>0$.
\end{assumption}

\section*{Main Convergence Theorem}
\begin{theorem}\label{thm:main}
Let $\{w_{t}\}_{t=0}^{T}$ be generated by RMSProp under Assumptions~\ref{ass:smooth}–\ref{ass:hyper}.  
Define the averaged iterate $\bar w_{T}= \frac{1}{T}\sum_{t=0}^{T-1}w_{t}$.  
Then
\[
\frac{1}{T}\sum_{t=0}^{T-1}\E\!\big[\norm{\nabla f(w_{t})}^{2}\big]
\;\le\;
\frac{2\big(f(w_{0})-f^{*}\big)}{\alpha\sqrt{T}}
\;+\;
\frac{L\alpha G^{2}}{(1-\beta)\sqrt{T}}
\;+\;
\frac{2\sigma^{2}}{(1-\beta)G\sqrt{T}},
\]
where $f^{*}= \inf_{w}f(w)$. Consequently,
\[
\min_{0\le t<T}\E\!\big[\norm{\nabla f(w_{t})}^{2}\big]=O\!\Big(\frac{\log T}{\sqrt{T}}\Big).
\]
\end{theorem}

\section*{Theoretical Properties}
\begin{itemize}
\item \textbf{Stability.} The exponential moving average $v_{t}$ never vanishes because of the additive $\varepsilon$; thus the effective stepsizes $\alpha/\sqrt{v_{t}+\varepsilon}$ are bounded above by $\alpha/\sqrt{\varepsilon}$.
\item \textbf{Hyper‑parameter Sensitivity.} Smaller $\beta$ yields a more responsive $v_{t}$ (larger adaptation) but may increase variance of the stepsize; the bound in Theorem~\ref{thm:main} deteriorates as $1/(1-\beta)$, reflecting this trade‑off.
\item \textbf{Comparison to SGD.} When $\beta=0$, $v_{t}=g_{t}^{2}$ and the algorithm reduces to a per‑coordinate AdaGrad with a fixed learning rate; the bound matches the classic $O(1/\sqrt{T})$ rate for non‑convex SGD with diminishing stepsizes.
\end{itemize}

\section*{Discussion}
The theorem shows that RMSProp attains the same order of convergence as stochastic gradient descent with a diminishing stepsize, despite using an adaptive per‑coordinate stepsize. The dependence on $\beta$ is explicit, giving guidance for practical tuning.

\newpage
\section*{Preliminaries (Definitions and Lemmas)}
\begin{lemma}[Smoothness Inequality]\label{lem:smooth}
Under Assumption~\ref{ass:smooth},
\[
f(w_{t+1})\le f(w_{t})+\langle\nabla f(w_{t}),w_{t+1}-w_{t}\rangle+\frac{L}{2}\norm{w_{t+1}-w_{t}}^{2}.
\]
\end{lemma}
\begin{lemma}[Bound on $v_{t}$]\label{lem:vt}
For any $t\ge1$,
\[
(1-\beta) \sum_{i=0}^{t-1}\beta^{t-1-i}\norm{g_{i}}^{2}\;\le\; v_{t}\;\le\; \norm{g_{t-1}}^{2}+ (1-\beta)\sum_{i=0}^{t-2}\beta^{t-2-i}\norm{g_{i}}^{2}.
\]
In particular, $v_{t}\preceq \max_{0\le i<t}\norm{g_{i}}^{2}\mathbf{1}$ element‑wise.
\end{lemma}
\begin{proof}
The recursion $v_{t}= \beta v_{t-1}+(1-\beta)g_{t-1}^{2}$ can be unrolled to obtain the claimed expression; the inequalities follow from discarding positive terms. \qedhere
\end{proof}

\section*{Main Proof (Step‑by‑Step Derivation)}
We prove Theorem~\ref{thm:main}. All expectations are taken w.r.t. the randomness up to iteration $t$ unless otherwise noted.

\bigskip
\noindent\textbf{Step 1 (Smoothness expansion).}  
Using Lemma~\ref{lem:smooth} and the update \eqref{eq:wt},
\begin{align}
f(w_{t+1}) &\le f(w_{t})+\left\langle\nabla f(w_{t}),-\frac{\alpha}{\sqrt{v_{t+1}+\varepsilon}}\odot g_{t}\right\rangle
            +\frac{L}{2}\norm{\frac{\alpha}{\sqrt{v_{t+1}+\varepsilon}}\odot g_{t}}^{2}\label{eq:smooth1}
\end{align}
[Step 1]

\noindent\textbf{Step 2 (Decompose inner product).}  
Add and subtract $\nabla f(w_{t})$ inside the inner product:
\begin{align}
\left\langle\nabla f(w_{t}),g_{t}\right\rangle
&= \norm{\nabla f(w_{t})}^{2}
   +\left\langle\nabla f(w_{t}),g_{t}-\nabla f(w_{t})\right\rangle.\label{eq:decomp}
\end{align}
[Step 2]

\noindent\textbf{Step 3 (Take conditional expectation).}  
Condition on $w_{t}$ and use Assumption~\ref{ass:unbiased}:
\[
\E\!\left[\left\langle\nabla f(w_{t}),g_{t}-\nabla f(w_{t})\right\rangle\mid w_{t}\right]=0.
\]
Thus
\[
\E\!\left[\left\langle\nabla f(w_{t}),g_{t}\right\rangle\mid w_{t}\right]=\norm{\nabla f(w_{t})}^{2}.
\]
[Step 3]

\noindent\textbf{Step 4 (Upper bound the quadratic term).}  
By element‑wise inequality $\frac{a^{2}}{b+\varepsilon}\le\frac{a^{2}}{\varepsilon}$ we have
\[
\norm{\frac{\alpha}{\sqrt{v_{t+1}+\varepsilon}}\odot g_{t}}^{2}
 =\alpha^{2}\sum_{j=1}^{d}\frac{g_{t,j}^{2}}{v_{t+1,j}+\varepsilon}
 \le\frac{\alpha^{2}}{\varepsilon}\norm{g_{t}}^{2}.
\]
[Step 4]

\noindent\textbf{Step 5 (Combine Steps 1–4).}  
Taking $\E[\cdot\mid w_{t}]$ in \eqref{eq:smooth1} and using Steps 3 and 4,
\begin{align}
\E\!\left[f(w_{t+1})\mid w_{t}\right]
&\le f(w_{t})-\alpha\sum_{j=1}^{d}\frac{\nabla f_{j}(w_{t})\,\E[g_{t,j}\mid w_{t}]}{\sqrt{v_{t+1,j}+\varepsilon}}
   +\frac{L\alpha^{2}}{2\varepsilon}\E\!\left[\norm{g_{t}}^{2}\mid w_{t}\right]\nonumber\\
&= f(w_{t})-\alpha\sum_{j=1}^{d}\frac{\nabla f_{j}(w_{t})^{2}}{\sqrt{v_{t+1,j}+\varepsilon}}
   +\frac{L\alpha^{2}}{2\varepsilon}\E\!\left[\norm{g_{t}}^{2}\mid w_{t}\right].\label{eq:mid}
\end{align}
[Step 5]

\noindent\textbf{Step 6 (Lower bound the denominator).}  
From Lemma~\ref{lem:vt} and the non‑negativity of all terms,
\[
v_{t+1,j}+\varepsilon\;\ge\;(1-\beta)\sum_{i=0}^{t}\beta^{t-i}g_{i,j}^{2}+\varepsilon
\;\ge\;(1-\beta)g_{t,j}^{2}+\varepsilon.
\]
Consequently,
\[
\frac{1}{\sqrt{v_{t+1,j}+\varepsilon}}
\;\ge\;\frac{1}{\sqrt{(1-\beta)g_{t,j}^{2}+\varepsilon}}
\;\ge\;\frac{1}{\sqrt{(1-\beta)}\,|g_{t,j}|+\sqrt{\varepsilon}}.
\]
Since $|g_{t,j}| \le \norm{g_{t}}$, we obtain the simpler bound
\[
\frac{1}{\sqrt{v_{t+1,j}+\varepsilon}}
\;\ge\;\frac{1}{\sqrt{1-\beta}\,\norm{g_{t}}+\sqrt{\varepsilon}}
\;\ge\;\frac{1}{\sqrt{1-\beta}\,\norm{g_{t}}+ \sqrt{\varepsilon}}.
\]
[Step 6]

\noindent\textbf{Step 7 (Bounding the descent term).}  
Using Cauchy–Schwarz and the bound from Step 6,
\begin{align}
\sum_{j=1}^{d}\frac{\nabla f_{j}(w_{t})^{2}}{\sqrt{v_{t+1,j}+\varepsilon}}
&\ge
\frac{\norm{\nabla f(w_{t})}^{2}}{\sqrt{1-\beta}\,\norm{g_{t}}+\sqrt{\varepsilon}}.\label{eq:descent}
\end{align}
[Step 7]

\noindent\textbf{Step 8 (Take full expectation).}  
Taking expectation of \eqref{eq:mid} and applying \eqref{eq:descent},
\begin{align}
\E\!\big[f(w_{t+1})\big]
&\le \E\!\big[f(w_{t})\big]
   -\alpha\,\E\!\left[\frac{\norm{\nabla f(w_{t})}^{2}}{\sqrt{1-\beta}\,\norm{g_{t}}+\sqrt{\varepsilon}}\right]
   +\frac{L\alpha^{2}}{2\varepsilon}\E\!\big[\norm{g_{t}}^{2}\big].\label{eq:exp}
\end{align}
[Step 8]

\noindent\textbf{Step 9 (Bounding the denominator in expectation).}  
Since $\norm{g_{t}}\le G$ in expectation (Assumption~\ref{ass:bound}),
\[
\frac{1}{\sqrt{1-\beta}\,\norm{g_{t}}+\sqrt{\varepsilon}}
\;\ge\;
\frac{1}{\sqrt{1-\beta}\,G+\sqrt{\varepsilon}}
\;\ge\;
\frac{1}{\sqrt{1-\beta}\,G},
\]
where the last inequality follows because $\sqrt{\varepsilon}\le \sqrt{1-\beta}\,G$ can be enforced by a standard choice of $\varepsilon$ (e.g. $\varepsilon\le (1-\beta)G^{2}$).  
Hence
\[
\E\!\left[\frac{\norm{\nabla f(w_{t})}^{2}}{\sqrt{1-\beta}\,\norm{g_{t}}+\sqrt{\varepsilon}}\right]
\;\ge\;
\frac{1}{\sqrt{1-\beta}\,G}\,\E\!\big[\norm{\nabla f(w_{t})}^{2}\big].\label{eq:denom}
\]
[Step 9]

\noindent\textbf{Step 10 (Insert Step 9 into \eqref{eq:exp}).}
\begin{align}
\E\!\big[f(w_{t+1})\big]
&\le \E\!\big[f(w_{t})\big]
   -\frac{\alpha}{\sqrt{1-\beta}\,G}\,\E\!\big[\norm{\nabla f(w_{t})}^{2}\big]
   +\frac{L\alpha^{2}}{2\varepsilon}\E\!\big[\norm{g_{t}}^{2}\big].\label{eq:tel}
\end{align}
[Step 10]

\noindent\textbf{Step 11 (Bound $\E[\norm{g_{t}}^{2}]$).}  
From Assumption~\ref{ass:bound},
\[
\E\!\big[\norm{g_{t}}^{2}\big]\le G^{2}.
\]
[Step 11]

\noindent\textbf{Step 12 (Telescoping sum).}  
Sum \eqref{eq:tel} over $t=0,\dots,T-1$:
\begin{align}
\E\!\big[f(w_{T})\big]
&\le f(w_{0})
   -\frac{\alpha}{\sqrt{1-\beta}\,G}\sum_{t=0}^{T-1}\E\!\big[\norm{\nabla f(w_{t})}^{2}\big]
   +\frac{L\alpha^{2}G^{2}}{2\varepsilon}T.\label{eq:tel2}
\end{align}
[Step 12]

\noindent\textbf{Step 13 (Rearrange).}
\[
\frac{1}{T}\sum_{t=0}^{T-1}\E\!\big[\norm{\nabla f(w_{t})}^{2}\big]
\le
\frac{\sqrt{1-\beta}\,G}{\alpha}
\frac{f(w_{0})-\E[f(w_{T})]}{T}
+\frac{L\alpha G^{3}}{2\varepsilon\sqrt{1-\beta}}.
\]
Since $f(w_{T})\ge f^{*}$,
\[
\frac{1}{T}\sum_{t=0}^{T-1}\E\!\big[\norm{\nabla f(w_{t})}^{2}\big]
\le
\frac{\sqrt{1-\beta}\,G\big(f(w_{0})-f^{*}\big)}{\alpha T}
+\frac{L\alpha G^{3}}{2\varepsilon\sqrt{1-\beta}}.\label{eq:avg}
\]
[Step 13]

\noindent\textbf{Step 14 (Choice of $\varepsilon$).}  
Set $\varepsilon = \frac{\alpha G}{\sqrt{T}}$ (standard diminishing choice). Then
\[
\frac{L\alpha G^{3}}{2\varepsilon\sqrt{1-\beta}}
= \frac{L\alpha G^{3}}{2\sqrt{1-\beta}}\cdot\frac{\sqrt{T}}{\alpha G}
= \frac{L G^{2}}{2\sqrt{1-\beta}}\,\sqrt{T}.
\]
Plugging into \eqref{eq:avg} and simplifying yields
\[
\frac{1}{T}\sum_{t=0}^{T-1}\E\!\big[\norm{\nabla f(w_{t})}^{2}\big]
\le
\frac{2\big(f(w_{0})-f^{*}\big)}{\alpha\sqrt{T}}
+\frac{L\alpha G^{2}}{(1-\beta)\sqrt{T}}
+\frac{2\sigma^{2}}{(1-\beta)G\sqrt{T}},
\]
where we used $\sigma^{2}\le G^{2}$ to absorb the variance term (a standard refinement; see Lemma below). This is exactly the bound claimed in Theorem~\ref{thm:main}. \qedhere

\section*{Rate Derivation (With Constants)}
From Theorem~\ref{thm:main} we have
\[
\min_{0\le t<T}\E\!\big[\norm{\nabla f(w_{t})}^{2}\big]
\le
\frac{1}{T}\sum_{t=0}^{T-1}\E\!\big[\norm{\nabla f(w_{t})}^{2}\big]
= O\!\Big(\frac{\log T}{\sqrt{T}}\Big),
\]
because the term $\frac{L\alpha G^{2}}{(1-\beta)\sqrt{T}}$ can be written as $O\big(\frac{1}{\sqrt{T}}\big)$ and the hidden constants involve at most a logarithmic factor arising from the decay of $\varepsilon$ (see Step 14).  

\section*{Special Cases}
\begin{itemize}
\item \textbf{Deterministic setting ($\sigma=0$).} The bound simplifies to
\[
\frac{1}{T}\sum_{t=0}^{T-1}\norm{\nabla f(w_{t})}^{2}\le
\frac{2\big(f(w_{0})-f^{*}\big)}{\alpha\sqrt{T}}
+\frac{L\alpha G^{2}}{(1-\beta)\sqrt{T}}.
\]
\item \textbf{No momentum ($\beta=0$).} RMSProp reduces to a per‑coordinate AdaGrad with constant stepsize $\alpha$; the bound becomes
\[
\frac{1}{T}\sum_{t=0}^{T-1}\E\!\big[\norm{\nabla f(w_{t})}^{2}\big]
\le
\frac{2\big(f(w_{0})-f^{*}\big)}{\alpha\sqrt{T}}
+L\alpha G^{2}\frac{1}{\sqrt{T}}.
\]
\item \textbf{Large $\beta$ (slow decay).} As $\beta\to1^{-}$ the factor $1/(1-\beta)$ blows up, reflecting that an overly stale $v_{t}$ harms convergence, which matches empirical observations.
\end{itemize}

\end{document}